\section{Светодиоды и таймеры}
\label{sec:leds}

\begin{infobox}[Цели главы]
\begin{itemize}
    \item Понять, как устроена светодиодная линейка Пионера и почему её адреса --- \texttt{0..28}.
    \item Задавать цвет через модуль \texttt{Ledbar} и модель RGB.
    \item Делать анимации на таймерах \texttt{Timer.new} и \texttt{Timer.callLater}, не останавливая программу.
    \item Реагировать на события автопилота в \texttt{callback(event)}.
    \item Строить индикацию как конечный автомат: состояние $\to$ событие $\to$ переход.
\end{itemize}
\end{infobox}

\subsection{Зачем дрону светодиоды}

Светодиоды --- первый видимый результат программы: дрон стоит на земле, а цвет меняется. Поэтому с них начинают. Но это ещё и рабочий инструмент: по цвету понятно, в каком состоянии программа --- готовится, летит, садится, сообщает об ошибке, --- даже когда терминала нет рядом. В главе~\ref{sec:missions} мы будем подсвечивать этапы миссии именно так.

В этой главе дрон не взлетает. Все примеры безопасно запускать и в симуляторе, и на настоящем Пионере со снятыми пропеллерами.

\subsection{Толковый словарь}
\begin{description}
  \item[\hypertarget{term:ledbar}{Ledbar}] модуль управления RGB-светодиодами: \verb|Ledbar.new(N)| создаёт объект, \verb|leds:set(i, r, g, b)| задаёт цвет.
  \item[\hypertarget{term:rgb}{RGB}] три компоненты цвета --- красная, зелёная, синяя --- числа от \verb|0.0| (выключено) до \verb|1.0| (максимум).
  \item[\hypertarget{term:index}{Индекс светодиода}] номер в линейке, от \verb|0| до \verb|N-1|.
  \item[\hypertarget{term:timer}{Таймер}] «будильник», который вызывает функцию через заданное время (\verb|Timer.callLater|) или периодически (\verb|Timer.new|).
  \item[\hypertarget{term:tick}{Тик}] один вызов функции таймера.
  \item[\hypertarget{term:event}{Событие}] сигнал от автопилота --- взлёт завершён, точка достигнута, батарея разряжена. Передаётся в \verb|callback(event)| как значение \verb|Ev.*|.
  \item[\hypertarget{term:state}{Состояние}] текущий режим логики программы, например \verb|"RED"| или \verb|"FLIGHT"|; хранится в переменной.
  \item[\hypertarget{term:automaton}{Конечный автомат}] программа, которая всегда находится ровно в одном состоянии и переходит в другое только по событию.
  \item[\hypertarget{term:blocking}{Блокирующий вызов}] функция, которая останавливает всю программу, пока не завершится, например \verb|sleep(t)|.
\end{description}
